% Reconstructed after the environment replacement; this restores the
% constructible-counting route independently of the quadratic witness proof.
\subsection{Constructible quotients and fixed-rank periodic counts}

\begin{lemma}[Proved: effective constructible quotients]
\label{D:effective-constructible-quotient}
Let \(k\) be algebraically closed, let \(T\) be a constructible set
of finite type over \(k\), and let \(\mathcal R\subset T\times T\)
be a constructible equivalence relation. There is a constructible
set \(Q\) and a constructible surjection \(c:T\to Q\) whose fibres
are exactly the equivalence classes. A constructible relation
on \(T\) inducing a partially defined single-valued operation on
classes descends to a partial function on \(Q\) with constructible
graph. The ground field need not be algebraic over a finite field.
\end{lemma}

\begin{proof}
The effective quotient assertion is
\cite[Theorem 2; constructive proof by Proposition 3.1, pp.~10--11]
{Beke2011FixedRankB}. Its scope is the constructible category over
an arbitrary field; the finite-field hypothesis in that paper's
zeta-function applications is not a hypothesis of the quotient
theorem. One concrete construction chooses, uniformly in each
nonempty constructible fibre, a nonempty finite subset of bounded
cardinality, with equal fibres receiving equal subsets. Induction
on the number of affine coordinates gives such a choice. Encoding
an unordered finite tuple by symmetric polynomial invariants
then distinguishes the disjoint equivalence classes.
For the last assertion, the image of the given relation under
\(c\times c\) is constructible by Chevalley's theorem. By the
assumed compatibility and uniqueness it is the graph of the
desired partial function.
\end{proof}

\begin{lemma}[Proved: periodic points of a partial constructible function]
\label{D:constructible-periodic-count}
Let \(Q\) be a constructible set over an algebraically closed field
\(k\), and let \(\Phi\) be a partially defined function on \(Q(k)\)
with constructible graph. Suppose that the set of points with a
defined \(f\)-step orbit returning to the initial point is finite
for each \(f\geq1\). There is a constant \(A\geq1\), depending on
this one graph, such that
\[
 \#\operatorname{Fix}(\Phi^f)\leq A^f\qquad(f\geq1).
 \tag{F1}
\]
No finite field of definition is required.
\end{lemma}

\begin{proof}
A constructible set has a finite disjoint partition by affine
locally closed sets. Adding finitely many coordinate tags embeds
it, in the constructible category, in one affine space
\(\mathbb A^a_k\). Write the graph as a union of \(L\) basic locally
closed pieces
\[
 \Gamma_i=\{(x,y):
 F_{i1}(x,y)=\cdots=F_{is_i}(x,y)=0,\quad g_i(x,y)\ne0\}.
\]
The finite products of the inequalities defining a piece may be
used for \(g_i\). Choose \(s\geq s_i\) and \(D\geq1\) bounding all
degrees, including those of \(u g_i-1\).
For a word \(\boldsymbol i=(i_0,\ldots,i_{f-1})\), impose the
closed affine equations
\[
 F_{i_jb}(x_j,x_{j+1})=0,\qquad
 u_jg_{i_j}(x_j,x_{j+1})=1,\qquad x_f=x_0.
\]
Use only the \(f\) variables \(x_0,\ldots,x_{f-1}\), so the last
equality is substituted rather than introducing more coordinates.
There are at most \((s+1)f\) equations of degree at most \(D\).
A solution is determined by \(x_0\): single-valuedness determines
each successive \(x_j\), and every \(u_j\) is uniquely determined.
The word scheme therefore has finitely many geometric points and
is zero-dimensional.

For clarity, the degree estimate needs no elimination of an
iterated graph. Intersect affine space successively with the
displayed equations. If an irreducible component is contained
in the next hypersurface, retain its degree. Otherwise the
hypersurface degree formula bounds the sum of the degrees of the
new components by \(D\) times its old degree. This formula follows
by projective closure and the exact sequence for multiplication
by a homogeneous equation, or directly from Hilbert polynomials.
Thus the word scheme has at most \(D^{(s+1)f}\) geometric points.
Summing over its \(L^f\) possible words gives (F1) with
\(A=LD^{s+1}\). A point may occur in several words, which only
increases this upper bound. Partiality causes no change:
a cycle solution uses only points where the next arrow is defined.
\end{proof}

\begin{corollary}[Proved: local fixed-rank representation growth]
\label{D:local-fixed-rank-representation-growth}
For a complete excellent normal local surface over an algebraically
closed field of characteristic \(p>0\), and for each fixed \(r\),
there is \(C_{R,r}\) such that the number of rank-\(r\)
\(\mathbf F_{p^f}\)-local systems on its puncture, up to isomorphism,
is at most \(p^{C_{R,r}f}\).
\end{corollary}

\begin{proof}
The uniform theorem \ref{D:uniform-local-representation-growth}
gives this with \(C_{R,r}=C_Rr^2\).
The earlier constructible proof uses the local atlas recorded
below, applies the effective quotient to ordinary reflexive-module
isomorphism, and uses reflexive Frobenius as its constructible
operation. The complete-local Lang lemma
\ref{D:local-complete-Lang} identifies periodic module classes
with unit structures. Finiteness for each fixed finite target
then allows \ref{D:constructible-periodic-count} to count cycles.
This explains the role of the quotient lemma in the independently
retained fixed-rank route; the stronger uniform proof does not
assume a complexity bound for that quotient.
\end{proof}

\begin{corollary}[Proved under stated hypotheses: simple-product quotients]
\label{D:bounded-matrix-simple-product}
Let \(G\) have the fixed-rank representation-class bounds of the
preceding corollary. Fix \(r_0\) and \(\epsilon>0\), and let
\(\mathcal S\) be a collection of pairwise nonisomorphic nonabelian
finite simple groups such that every \(S\in\mathcal S\) admits
a faithful representation
\[
 S\hookrightarrow\operatorname{GL}_{r_0}(\mathbf F_{p^{f_S}}),
 \qquad |S|\geq p^{\epsilon f_S}.
\]
Every quotient of \(G\) which is a product of groups in
\(\mathcal S\) is topologically finitely generated.
This includes any family for which these fixed matrix-size and
field-size inequalities have been verified.
\end{corollary}

\begin{proof}
Write such a quotient as \(\prod_{S\in\mathcal S}S^{m_S}\).
Its coordinate epimorphisms have distinct kernels.
Composing them with the fixed faithful representation of \(S\)
gives distinct conjugacy classes of representations: conjugate
representations have equal kernels. Hence
\[
 m_S\leq p^{C_{r_0}f_S}\leq |S|^{C_{r_0}/\epsilon}.
 \tag{F2}
\]
We recall the precise finite-group input used to turn (F2) into
generation. In the proof of Proposition 2.6(ii), printed p.~12,
of \cite{LucchiniMarionTracey2020D}, used in
\ref{D:cohomology-crown-criterion}, the number of
\(\operatorname{Aut}(S)\)-orbits of generating \(s\)-tuples is at
least
\[
 \frac{|S|^{s-1}}{2|\operatorname{Out}(S)|}\qquad(s\geq2),
\]
and the same passage gives two-generation of \(S\); p.~13 records
\(|\operatorname{Out}(S)|\leq\log |S|\).
Thus one fixed sufficiently large \(s\) supplies at least \(m_S\)
distinct generating-tuple orbits for every \(S\).
Choose those tuples at the coordinate copies.
Goursat's lemma shows that they generate every finite subproduct:
for equal simple factors the distinct automorphism orbits exclude
diagonal relations, and nonisomorphic simple factors have no
common simple quotient. Their \(s\) assembled elements therefore
generate a dense subgroup of the entire product.
This proves the assertion without claiming anything about
unbounded matrix dimension or general monolithic crowns.
\end{proof}

